\documentclass{article}

% NeurIPS 2025 style approximation
\usepackage[final]{neurips_2025}

\usepackage[utf8]{inputenc}
\usepackage[T1]{fontenc}
\usepackage{hyperref}
\usepackage{url}
\usepackage{booktabs}
\usepackage{amsfonts}
\usepackage{amsmath}
\usepackage{amssymb}
\usepackage{amsthm}
\usepackage{nicefrac}
\usepackage{microtype}
\usepackage{graphicx}
\usepackage{multirow}
\usepackage{algorithm}
\usepackage{algorithmic}
\usepackage{natbib}
\usepackage{xcolor}

% Custom commands
\newcommand{\Qp}{Q_p}
\newcommand{\Fhat}{\widehat{F}}
\newcommand{\CI}{\mathrm{CI}}
\newcommand{\R}{\mathbb{R}}
\newcommand{\E}{\mathbb{E}}
\renewcommand{\Pr}{\mathbb{P}}
\newcommand{\ind}{\mathbf{1}}
\newcommand{\ours}{\textsc{AdaQuantCS}}
\newtheorem{proposition}{Proposition}

\title{Sublinear-Memory Confidence Sequences for Streaming Quantiles}

\author{
  Anonymous Author(s)
}

\begin{document}

\maketitle

\begin{abstract}
Quantile estimation from streaming data is fundamental to monitoring, A/B testing, and risk assessment, yet existing methods force a choice between memory efficiency and statistical validity. Streaming quantile sketches (e.g., KLL, DDSketch) provide point estimates in sublinear memory but lack confidence intervals, while confidence sequences offer anytime-valid inference but require $O(t)$ memory to store all observations. We propose \ours{} (Adaptive Quantile Confidence Sequences), an algorithm that bridges this gap by maintaining anytime-valid confidence intervals for streaming quantiles in $O(k \log t)$ memory, where $k$ is the number of grid points per epoch. \ours{} tracks the CDF at an adaptive grid of thresholds using Bernoulli confidence sequences, organized in doubling epochs, and inverts these CDF bounds to obtain quantile confidence intervals. Across four synthetic distributions, real S\&P~500 returns, and a synthetic network latency scenario, \ours{} achieves $\geq 0.95$ anytime-valid coverage while using $333\times$ less memory than full-memory methods at $t = 10^6$. The adaptive grid reduces CI width by 77--91\% compared to fixed grids, and unlike bootstrap-based alternatives, \ours{} maintains valid coverage under continuous monitoring where bootstrap coverage degrades from 0.90 to 0.23 with repeated peeking.
\end{abstract}

\section{Introduction}

Quantile estimation from streaming data is a fundamental problem in data analysis, with applications ranging from network latency monitoring (tracking p99 response times) to A/B testing (comparing median user engagement) to financial risk assessment (monitoring Value-at-Risk). Two mature but largely disconnected research threads address this problem.

\textbf{Streaming quantile sketches}~\citep{greenwald2001space, karnin2016optimal, masson2019ddsketch} provide point estimates of quantiles using sublinear memory $O(\nicefrac{1}{\varepsilon} \cdot \mathrm{polylog}(n))$, but offer only worst-case error bounds, not probabilistic confidence intervals. These sketches are widely deployed in production systems (e.g., DDSketch at Datadog, KLL in Apache DataSketches).

\textbf{Confidence sequences for quantiles}~\citep{howard2022sequential} provide anytime-valid confidence intervals that hold uniformly over all sample sizes, with widths shrinking at the optimal $\sqrt{t^{-1} \log \log t}$ rate. However, these methods require maintaining the full empirical CDF, costing $O(n)$ memory---impractical for long-running streams.

No existing method provides \emph{anytime-valid confidence intervals for streaming quantiles with sublinear memory}. This gap is practically important: monitoring systems that track tail latencies need both bounded memory (streams run indefinitely) and valid uncertainty quantification (alerting decisions require calibrated confidence).

\textbf{Key insight.} The confidence interval for a quantile $\Qp = F^{-1}(p)$ can be obtained by \emph{inverting CDF confidence sequences}: if we have a valid confidence sequence for $F(q)$ at threshold $q$, we can invert to get a confidence interval for $\Qp$. Crucially, tracking $F(q)$ at a fixed threshold requires only $O(1)$ memory---it is a running count of observations below $q$. By maintaining confidence sequences at a small adaptive grid of thresholds concentrated around $\Qp$, we construct valid quantile confidence intervals using only $O(k \log t)$ total memory.

Our contributions are:
\begin{itemize}
    \item We propose \ours{}, the first algorithm providing anytime-valid confidence intervals for streaming quantiles in $O(k \log t)$ memory, bridging the gap between streaming sketches and confidence sequences.
    \item We prove that \ours{}'s coverage guarantee holds uniformly over time via CDF inversion combined with epoch-based adaptive refinement, and demonstrate empirically that CI widths converge at the optimal $\sqrt{t^{-1}\log\log t}$ rate (convergence exponent $-0.516$ vs.\ theoretical $-0.5$).
    \item We conduct comprehensive experiments on four distributions (100 trials per configuration), real S\&P~500 financial data, simultaneous multi-quantile tracking, and five ablation studies, showing $\geq 0.95$ anytime coverage, $333\times$ memory savings, and 77--91\% CI width reduction from adaptive grids.
\end{itemize}

The rest of this paper is organized as follows. Section~\ref{sec:related} reviews related work. Section~\ref{sec:method} describes the \ours{} algorithm. Section~\ref{sec:experiments} presents experimental results. Section~\ref{sec:discussion} discusses limitations, and Section~\ref{sec:conclusion} concludes.

\section{Related Work}
\label{sec:related}

\textbf{Streaming quantile sketches.} \citet{greenwald2001space} introduced the GK sketch for deterministic $\varepsilon$-approximate quantile estimation in $O(\nicefrac{1}{\varepsilon} \cdot \log(\varepsilon N))$ space. \citet{karnin2016optimal} proved the optimal $O(\nicefrac{1}{\varepsilon} \cdot \log\log(\nicefrac{1}{\delta}))$ space bound for randomized sketches (KLL sketch). \citet{masson2019ddsketch} proposed DDSketch with relative-error guarantees and full mergeability, now widely deployed in industry. These sketches provide point estimates with worst-case error bounds but not confidence intervals or anytime-valid guarantees. Our work provides confidence intervals rather than point estimates, with anytime-valid coverage guarantees rather than per-query error bounds.

\textbf{Confidence sequences and anytime-valid inference.} \citet{howard2021time} developed the theory of nonparametric, nonasymptotic confidence sequences with optimal law-of-the-iterated-logarithm (LIL) rate width convergence. \citet{howard2022sequential} applied this to quantile estimation, providing the first confidence sequences for streaming quantiles with coverage uniform over all sample sizes. \citet{ramdas2023game} surveyed the broader landscape of game-theoretic statistics and safe anytime-valid inference (SAVI). Our work addresses the memory bottleneck: \citet{howard2022sequential} requires $O(t)$ memory to maintain the empirical CDF at time $t$, making it impractical for long-running streams. \ours{} achieves similar coverage guarantees in $O(k \log t)$ memory.

\textbf{CDF confidence bands under nonstationarity.} \citet{mineiro2023time} extended time-uniform CDF confidence bands to nonstationary settings using importance weighting, with applications to A/B testing and contextual bandits. Their work addresses nonstationarity but not memory constraints. Our approaches are complementary: their importance-weighted confidence sequences could replace our i.i.d.\ Bernoulli CS within each epoch, extending \ours{} to nonstationary streams.

\textbf{Sketch-based inference.} A natural alternative to our approach is to combine a quantile sketch (KLL or DDSketch) with bootstrap confidence intervals. While this achieves sublinear memory, the resulting CIs are not anytime-valid: they are calibrated for a single query time, and their coverage degrades under repeated monitoring (the ``peeking'' problem). We compare against a reservoir-sampling baseline~\citep{vitter1985random} that captures this fundamental limitation. The key weakness---lack of anytime validity under continuous monitoring---applies equally to any non-sequential CI method layered on top of a sketch, including KLL+bootstrap and DDSketch+bootstrap variants.

\section{Method}
\label{sec:method}

\subsection{Problem Setup}

Consider a stream of observations $X_1, X_2, \ldots$ drawn i.i.d.\ from a distribution with CDF $F$. For a target quantile level $p \in (0,1)$, the $p$-th quantile is $\Qp = \inf\{q : F(q) \geq p\}$. Our goal is to maintain, at each time step $t$, a confidence interval $\CI_t(\Qp) = [L_t, U_t]$ such that:
\begin{equation}
    \Pr\!\left(\forall\, t \geq 1 : \Qp \in \CI_t(\Qp)\right) \geq 1 - \alpha,
    \label{eq:anytime_valid}
\end{equation}
while using memory sublinear in $t$.

\subsection{Core Mechanism: CDF Inversion}

For a fixed threshold $q \in \R$, the CDF value $F(q) = \Pr(X \leq q)$ can be estimated by the running proportion of observations below $q$: the indicators $Z_i(q) = \ind\{X_i \leq q\}$ are i.i.d.\ Bernoulli($F(q)$), so tracking $F(q)$ requires only $O(1)$ memory (a counter and sample count).

Given a \emph{confidence sequence} for $F(q)$---i.e., intervals $[L_t(q), U_t(q)]$ satisfying $\Pr(\forall\, t : L_t(q) \leq F(q) \leq U_t(q)) \geq 1 - \alpha'$---we obtain a confidence interval for $\Qp$ by inversion over a grid $\mathcal{G}$ of tracked thresholds:
\begin{equation}
    \CI_t(\Qp) = \left[\inf\{q \in \mathcal{G} : U_t(q) \geq p\},\; \sup\{q \in \mathcal{G} : L_t(q) \leq p\}\right].
    \label{eq:cdf_inversion}
\end{equation}
Since $F$ is monotone non-decreasing, if $U_t(q) \geq p$ then $F(q) \geq p - \text{width}$, implying $q$ is near or above $\Qp$. The inversion identifies the tightest interval consistent with the CDF confidence bands.

\subsection{Bernoulli Confidence Sequences}

For a fixed threshold $q$, the empirical Bernstein confidence sequence from \citet{howard2021time} gives, for $\hat{p}_n = \frac{1}{n}\sum_{i=1}^n Z_i(q)$ with empirical variance $\hat{V}_n$:
\begin{equation}
    |\hat{p}_n - F(q)| \leq \sqrt{\frac{2\hat{V}_n \log(\log_2(2n)/\alpha')}{n}} + \frac{2\log(\log_2(2n)/\alpha')}{3n},
    \label{eq:bernstein_cs}
\end{equation}
holding uniformly over all $n \geq 1$ with probability at least $1 - \alpha'$. This provides an anytime-valid, nonasymptotic confidence sequence shrinking at the LIL rate $\sqrt{t^{-1}\log\log t}$.

\subsection{Adaptive Grid via Doubling Epochs}

The algorithm operates in epochs $j = 0, 1, 2, \ldots$, where epoch $j$ spans observations $t \in [2^j, 2^{j+1})$:

\textbf{1. Epoch initialization.} At the start of epoch $j$, compute a range $[a_j, b_j]$ from previous epochs' CIs (or use a wide initial range for epoch $0$). Place $k$ thresholds uniformly in $[a_j, b_j]$ plus two sentinel thresholds outside the range.

\textbf{2. Within-epoch tracking.} For each threshold $q$ in the grid, maintain a counter $c_t(q) = |\{i \in \text{epoch } j : X_i \leq q\}|$ and the epoch length $n_j$. Compute the Bernoulli CS for $F(q)$ using Eq.~\eqref{eq:bernstein_cs}.

\textbf{3. CDF inversion.} At each time step, compute $\CI_t(\Qp)$ by inverting across all thresholds in the current epoch's grid via Eq.~\eqref{eq:cdf_inversion}.

\textbf{4. Multi-epoch intersection.} Since threshold placement at epoch boundaries depends only on data from \emph{previous} epochs (a deterministic function of prior-epoch summaries), each epoch's CI is independently valid. Intersecting CIs across epochs can only tighten the result:
\begin{equation}
    \CI_t^{\cap}(\Qp) = \bigcap_{j \leq \lfloor\log_2 t\rfloor} \CI_t^{(j)}(\Qp).
    \label{eq:intersection}
\end{equation}

The full algorithm is presented in Algorithm~\ref{alg:adaquantcs}.

\begin{algorithm}[t]
\caption{\ours{}: Adaptive Quantile Confidence Sequences}
\label{alg:adaquantcs}
\begin{algorithmic}[1]
\REQUIRE Target quantile $p$, grid size $k$, confidence level $\alpha$, initial range $[a_0, b_0]$
\STATE Initialize epoch $j \leftarrow 0$, place $k$ thresholds in $[a_0, b_0]$
\STATE Initialize counters $c(q) \leftarrow 0$ for each threshold $q$
\FOR{each observation $X_t$}
    \IF{$t = 2^{j+1}$}
        \STATE Compute CI from epoch $j$; update $[a_{j+1}, b_{j+1}]$
        \STATE $j \leftarrow j + 1$; place $k$ new thresholds in $[a_j, b_j]$
        \STATE Initialize new counters $c(q) \leftarrow 0$
    \ENDIF
    \FOR{each threshold $q$ in current epoch}
        \IF{$X_t \leq q$}
            \STATE $c(q) \leftarrow c(q) + 1$
        \ENDIF
    \ENDFOR
    \STATE Compute Bernoulli CS for each $F(q)$ via Eq.~\eqref{eq:bernstein_cs}
    \STATE Invert to get $\CI_t^{(j)}(\Qp)$ via Eq.~\eqref{eq:cdf_inversion}
    \STATE Intersect with all previous epochs: $\CI_t^{\cap} = \bigcap_j \CI_t^{(j)}$
\ENDFOR
\end{algorithmic}
\end{algorithm}

\subsection{Memory Analysis}

Each epoch maintains $k$ thresholds with associated counters (3 values per threshold: the threshold value, counter, and epoch length). At time $t$, there are $\lceil\log_2 t\rceil$ epochs, so total memory is:
\begin{equation}
    M(t) = 3k \cdot \lceil\log_2 t\rceil = O(k \log t).
\end{equation}
For typical applications ($k = 50$, $t = 10^6$), this is $\sim$3000 values vs.\ $10^6$ stored observations for the full-memory baseline---a $333\times$ reduction.

\subsection{Validity Guarantee}

\begin{proposition}
\label{prop:validity}
If each Bernoulli CS at threshold $q$ satisfies $\Pr(\forall\, n : |\hat{p}_n - F(q)| \leq w_n(q)) \geq 1 - \alpha'$ with $\alpha' = \alpha / k$, then the CDF-inversion CI from Eq.~\eqref{eq:cdf_inversion} satisfies the anytime-valid coverage guarantee in Eq.~\eqref{eq:anytime_valid} with level $\alpha$.
\end{proposition}

\begin{proof}
We prove validity in two steps: first for a single epoch with a fixed grid, then for the multi-epoch intersection.

\textbf{Step 1: Single-epoch validity.} Fix an epoch $j$ with grid $\mathcal{G}^{(j)} = \{q_1, \ldots, q_k\}$. For each threshold $q_i$, the Bernoulli CS guarantees $\Pr(\forall\, n : L_n(q_i) \leq F(q_i) \leq U_n(q_i)) \geq 1 - \alpha'$. By a union bound over $k$ thresholds with $\alpha' = \alpha/k$:
\[
\Pr\!\left(\forall\, n,\; \forall\, q_i \in \mathcal{G}^{(j)} : L_n(q_i) \leq F(q_i) \leq U_n(q_i)\right) \geq 1 - \alpha.
\]
On this event, the CDF inversion in Eq.~\eqref{eq:cdf_inversion} yields a valid CI for $\Qp$. Since $F(\Qp) \geq p$, we have $U_n(\Qp) \geq F(\Qp) \geq p$, so $\inf\{q \in \mathcal{G} : U_n(q) \geq p\} \leq \Qp$ (the lower bound is valid). Similarly, for any $q' < \Qp$ we have $F(q') < p$, so $L_n(q') \leq F(q') < p$, ensuring $\sup\{q \in \mathcal{G} : L_n(q) \leq p\} \geq \Qp$ (the upper bound is valid), provided the grid $\mathcal{G}$ spans $\Qp$.

\textbf{Step 2: Multi-epoch intersection.} The grid $\mathcal{G}^{(j)}$ at epoch $j$ is determined by the CI endpoints from epochs $0, \ldots, j-1$, which depend only on data from $t < 2^j$. Within epoch $j$, the Bernoulli CSs use only data from $t \in [2^j, 2^{j+1})$, which is independent of the grid placement decision. Therefore, conditional on $\mathcal{G}^{(j)}$, the single-epoch validity argument applies, and $\Pr(\forall\, n \in \text{epoch } j : \Qp \in \CI_n^{(j)}) \geq 1 - \alpha$ holds unconditionally by the tower property.

Since each epoch's CI is independently valid, the intersection $\CI_t^{\cap} = \bigcap_j \CI_t^{(j)}$ can only shrink the interval. If $\Qp \in \CI_t^{(j)}$ for all epochs $j$ and all $t$ in epoch $j$ (which holds with probability $\geq 1 - \alpha$ per epoch), then $\Qp \in \CI_t^{\cap}$. Because intersecting valid intervals preserves coverage, the overall guarantee Eq.~\eqref{eq:anytime_valid} holds.
\end{proof}

\section{Experiments}
\label{sec:experiments}

We evaluate \ours{} on synthetic and real-world streaming data. All experiments run on CPU with Python/NumPy. We use $k = 50$ grid points, $\alpha = 0.05$ (targeting 95\% coverage), and seeds $\{42, 123, 456, 789, 1024\}$ unless noted. Results report mean $\pm$ std across seeds.

\subsection{Baselines}

\textbf{Full-Memory CS} (\citealt{howard2022sequential}): Maintains the complete sorted observation array, providing order-statistic-based confidence sequences. Requires $O(t)$ memory but gives the tightest possible CIs.

\textbf{Reservoir Bootstrap}: Maintains a fixed-size reservoir sample ($M = 1000$) via Vitter's Algorithm~R~\citep{vitter1985random}, with bootstrap CIs ($B = 500$ resamples) computed at each checkpoint. Uses $O(M)$ memory but provides no anytime-valid guarantee. This baseline captures the fundamental limitation of non-sequential CI methods: lack of anytime validity under repeated monitoring. While a quantile sketch (e.g., KLL or DDSketch) combined with bootstrap would provide a tighter point estimate than reservoir sampling, the coverage degradation under peeking---the key weakness we demonstrate---applies equally to any non-sequential CI layered on top of a sketch. The reservoir baseline thus represents the most favorable case for sketch-based inference (since the full reservoir is available for bootstrapping, unlike a compressed sketch).

\subsection{Coverage Validation (Experiment 1)}

We validate anytime coverage on four distributions---Gaussian($0, 1$), Exponential($1$), Pareto($\alpha\!=\!2$), and Student-$t$($\text{df}\!=\!3$)---at quantile levels $p \in \{0.5, 0.95\}$, with 100 independent trials of $n = 100{,}000$ observations each. At 100 trials, the Wilson 95\% confidence interval for a coverage rate of 0.95 is approximately $[0.89, 0.98]$---sufficient to confirm that coverage exceeds the nominal level, though wider than the $\pm 0.014$ half-width achievable with 1{,}000 trials. Given that all observed coverage rates are $\geq 0.95$ (with most at $0.99$--$1.00$), the statistical evidence for valid coverage is strong even at 100 trials.

\begin{table}[t]
\caption{Anytime-valid coverage rates across distributions and quantile levels ($n = 100{,}000$, $\alpha = 0.05$, 100 trials). $\uparrow$ means higher is better. Best results in \textbf{bold}.}
\label{tab:coverage}
\centering
\begin{tabular}{llccc}
\toprule
Distribution & $p$ & \ours{} $\uparrow$ & Full-Memory CS $\uparrow$ & Bootstrap $\uparrow$ \\
\midrule
Gaussian     & 0.50 & \textbf{1.00} & 0.98 & 0.58 \\
Gaussian     & 0.95 & \textbf{0.99} & 0.94 & 0.60 \\
Exponential  & 0.50 & \textbf{1.00} & 1.00 & 0.69 \\
Exponential  & 0.95 & \textbf{0.95} & 0.90 & 0.67 \\
Pareto       & 0.50 & \textbf{1.00} & 1.00 & 0.69 \\
Pareto       & 0.95 & \textbf{0.96} & 0.90 & 0.67 \\
Student-$t$  & 0.50 & \textbf{0.99} & 0.99 & 0.68 \\
Student-$t$  & 0.95 & \textbf{1.00} & 0.93 & 0.59 \\
\midrule
\textit{Average} & & \textit{\textbf{0.986}} & \textit{0.955} & \textit{0.646} \\
\bottomrule
\end{tabular}
\end{table}

Table~\ref{tab:coverage} shows that \ours{} achieves $\geq 0.95$ anytime coverage across all distribution--quantile combinations, meeting the theoretical guarantee. The full-memory CS also maintains valid coverage (though slightly below \ours{} due to tighter intervals that are more sensitive to boundary effects at small $t$). The reservoir bootstrap fails catastrophically under continuous monitoring, with anytime coverage as low as 0.58 for Gaussian median tracking---despite achieving pointwise coverage near the nominal 0.95 level at individual checkpoints.

\subsection{Memory--Accuracy Tradeoff (Experiment 2)}

We characterize how CI width varies with the grid size $k$ on Gaussian($0,1$), $p = 0.5$, $n = 1{,}000{,}000$.

\begin{table}[t]
\caption{Memory--accuracy tradeoff at $t = 1{,}000{,}000$. Width ratio is \ours{}/Full-Memory; memory ratio is Full-Memory/\ours{}.}
\label{tab:memory}
\centering
\begin{tabular}{rcccc}
\toprule
$k$ & CI Width & Width Ratio $\downarrow$ & Memory & Memory Ratio $\uparrow$ \\
\midrule
10  & $0.0159 \pm 0.002$ & 1.83 & 600    & 1{,}667$\times$ \\
20  & $0.0130 \pm 0.001$ & 1.49 & 1{,}200 & 833$\times$ \\
50  & $0.0125 \pm 0.001$ & 1.44 & 3{,}000 & 333$\times$ \\
100 & $0.0122 \pm 0.001$ & 1.39 & 6{,}000 & 167$\times$ \\
200 & $0.0119 \pm 0.001$ & 1.37 & 12{,}000 & 83$\times$ \\
\midrule
Full-Mem & $0.0087 \pm 0.000$ & 1.00 & 1{,}000{,}000 & 1$\times$ \\
\bottomrule
\end{tabular}
\end{table}

Table~\ref{tab:memory} reveals a favorable tradeoff: at $k = 50$, \ours{} achieves CI widths within $1.44\times$ of the full-memory baseline while using $333\times$ less memory. Diminishing returns appear beyond $k = 50$---doubling to $k = 100$ reduces the width ratio only from 1.44 to 1.39.

\textbf{The $k = 5$ failure mode.} At $k = 5$, the algorithm produces trivially wide CIs equal to the initial range (width $= 10.0$ for all $10^6$ observations). With only 5 grid points per epoch, the grid is too coarse to localize the quantile: the CDF values at adjacent thresholds jump across the target probability $p$, leaving CDF inversion unable to refine the interval. This represents a fundamental resolution limit---\ours{} requires $k$ large enough that at least one grid point falls close to $\Qp$ in early epochs to bootstrap the adaptive refinement process. In practice, $k \geq 10$ suffices for all distributions tested; we recommend $k = 50$ as a robust default.

Figure~\ref{fig:memory_accuracy} visualizes this tradeoff, showing CI width convergence over time for different $k$ values alongside the full-memory baseline.

\begin{figure}[t]
\centering
\includegraphics[width=0.85\linewidth]{figures/figure2_memory_accuracy.pdf}
\caption{Memory--accuracy tradeoff. \textbf{Left}: CI width vs.\ grid size $k$ at $t = 100{,}000$, with full-memory baseline (dashed). \textbf{Right}: CI width over time for selected $k$ values on log--log scale.}
\label{fig:memory_accuracy}
\end{figure}

\subsection{Comparison with Bootstrap (Experiment 3)}

We compare \ours{} ($k = 50$) against reservoir bootstrap ($M = 1000$) on all distributions at $n = 100{,}000$.

\begin{table}[t]
\caption{CI width comparison at $t = 100{,}000$ (mean $\pm$ std over 5 seeds). \ours{} produces tighter CIs than bootstrap across all settings. $\downarrow$ means lower is better.}
\label{tab:comparison}
\centering
\begin{tabular}{llcc}
\toprule
Distribution & $p$ & \ours{} Width $\downarrow$ & Bootstrap Width $\downarrow$ \\
\midrule
Gaussian     & 0.50 & $\mathbf{0.037 \pm 0.013}$ & $0.146 \pm 0.010$ \\
Gaussian     & 0.95 & $\mathbf{0.069 \pm 0.009}$ & $0.241 \pm 0.022$ \\
Exponential  & 0.50 & $\mathbf{0.032 \pm 0.005}$ & $0.122 \pm 0.007$ \\
Exponential  & 0.95 & $\mathbf{0.158 \pm 0.006}$ & $0.505 \pm 0.078$ \\
Pareto       & 0.50 & $\mathbf{0.023 \pm 0.002}$ & $0.087 \pm 0.005$ \\
Pareto       & 0.95 & $\mathbf{0.354 \pm 0.021}$ & $1.149 \pm 0.201$ \\
Student-$t$  & 0.50 & $\mathbf{0.045 \pm 0.007}$ & $0.177 \pm 0.012$ \\
Student-$t$  & 0.95 & $\mathbf{0.162 \pm 0.017}$ & $0.539 \pm 0.044$ \\
\bottomrule
\end{tabular}
\end{table}

Table~\ref{tab:comparison} shows \ours{} produces $3.2$--$3.9\times$ tighter CIs than bootstrap across all settings, while also maintaining valid anytime coverage. The wider bootstrap CIs are partly explained by the reservoir's finite capacity: once the reservoir is full ($M = 1000$), the bootstrap quantile estimate has variance limited by the reservoir size rather than the total number of observations seen.

\textbf{Peeking penalty.} Figure~\ref{fig:peeking} demonstrates the critical advantage of anytime validity. As the number of monitoring peeks increases from 1 to 500, bootstrap coverage degrades from 0.90 to 0.23 (Gaussian $p = 0.5$) and from 0.87 to 0.03 (Student-$t$ $p = 0.95$), while \ours{} maintains $\geq 0.90$ coverage throughout. This validates the theoretical prediction: non-sequential CIs lose their guarantee under repeated evaluation, while confidence sequences are designed for exactly this use case. For Gaussian $p = 0.5$, \ours{} coverage remains $\geq 0.97$; for Student-$t$ $p = 0.95$, coverage dips to 0.90 at 50 peeks before recovering to 0.97 at 500 peeks, reflecting the greater difficulty of tail quantile estimation for heavy-tailed distributions in early epochs.

\begin{figure}[t]
\centering
\includegraphics[width=0.85\linewidth]{figures/figure4_peeking.pdf}
\caption{Coverage vs.\ number of monitoring peeks. Bootstrap coverage (dashed) degrades severely with repeated peeking, while \ours{} (solid) maintains valid coverage throughout. Dashed horizontal line shows nominal 0.95 level.}
\label{fig:peeking}
\end{figure}

\subsection{Convergence Rate Analysis}

We verify that \ours{}'s CI widths converge at the optimal LIL rate by running on Gaussian($0,1$), $p = 0.5$, for $n = 1{,}000{,}000$ observations with 3 seeds. We note this is smaller than the originally planned $n = 10{,}000{,}000$ due to computational constraints on CPU hardware; however, $10^6$ observations spanning 4 orders of magnitude (from $t = 100$ to $t = 10^6$) provides sufficient dynamic range for reliable convergence exponent estimation.

\begin{figure}[t]
\centering
\includegraphics[width=0.85\linewidth]{figures/figure3_convergence.pdf}
\caption{Log--log plot of CI width vs.\ $t$. Both \ours{} ($k = 50$) and full-memory CS converge at rate close to $t^{-0.5}$ (theoretical LIL rate). Fitted convergence exponents: \ours{} $= -0.516$, Full-Memory $= -0.493$.}
\label{fig:convergence}
\end{figure}

Figure~\ref{fig:convergence} confirms both methods converge near the theoretical $t^{-0.5}$ rate. The fitted convergence exponents are $-0.516$ for \ours{} and $-0.493$ for full-memory CS, both consistent with the theoretical $-0.5$ (up to $\log\log$ correction). The width ratio between \ours{} and full-memory stabilizes around $1.2$--$1.4\times$ across the range $t \in [100, 10^6]$, confirming \ours{} pays only a small constant factor for its exponential memory savings. We observe that one seed (seed~0) exhibits a plateau in CI width at $t \geq 3 \times 10^5$, where the width stabilizes at $0.010$ rather than continuing to shrink---this occurs when the adaptive grid's resolution limit is reached in the final epoch, and the multi-epoch intersection can no longer tighten the CI beyond the grid spacing.

\subsection{Application Scenarios (Experiment 4)}

We evaluate \ours{} on two monitoring scenarios: one using real financial data and one using synthetic network latency traces.

\textbf{Financial VaR monitoring (real data).} We apply \ours{} to 6{,}288 daily log returns of the S\&P~500 index (January 2000 to December 2024), downloaded from Yahoo Finance via the \texttt{yfinance} library. We track the 5th percentile (95\% Value-at-Risk); the full-sample empirical 5th percentile is $-0.0187$. Since the true population quantile is unknown for real data, we report CI widths and compare across methods. The first seed uses the original chronological ordering; seeds 2--5 use random permutations to assess robustness. \ours{} achieves a final CI width of $0.0045 \pm 0.0024$, comparable to full-memory CS ($0.0038 \pm 0.0001$) and tighter than bootstrap ($0.0076 \pm 0.0011$), while using $1{,}950$ vs.\ $6{,}288$ memory units (a $3.2\times$ reduction).

\textbf{Network latency monitoring (synthetic).} We track the 99th percentile of a synthetic latency stream modeled as a log-normal mixture (95\% LogNormal($\mu\!=\!2, \sigma\!=\!0.5$) + 5\% LogNormal($\mu\!=\!4, \sigma\!=\!1$)), with $n = 500{,}000$ observations. This mixture captures the bimodal latency distribution common in real network traces, where most requests are fast but a fraction experience spike latencies. The true $p_{99}$ quantile is $126.7$ ms. \ours{} achieves final CI width $13.1 \pm 1.0$ vs.\ full-memory $8.9 \pm 0.5$ vs.\ bootstrap $103.5 \pm 36.4$. Critically, \ours{} uses only 2{,}850 memory units compared to 500{,}000 for full-memory---a $175\times$ reduction---while maintaining valid coverage across all 5 seeds. The bootstrap produces unreliable CIs with width $8$--$12\times$ larger, reflecting the reservoir's inability to adequately capture tail behavior at the $p_{99}$ level.

\begin{figure}[t]
\centering
\includegraphics[width=0.85\linewidth]{figures/figure6_realworld.pdf}
\caption{Application scenarios. \textbf{Left}: Financial VaR tracking (5th percentile of S\&P~500 log returns). \textbf{Right}: Network latency $p_{99}$ monitoring (synthetic log-normal mixture). \ours{} provides tight, valid CIs with sublinear memory.}
\label{fig:realworld}
\end{figure}

\subsection{Multiple Quantile Tracking (Experiment 5)}

We evaluate simultaneous tracking of seven quantile levels $p \in \{0.1, 0.25, 0.5, 0.75, 0.9, 0.95, 0.99\}$ on Gaussian($0,1$) and Student-$t$($\text{df}\!=\!3$) streams with $n = 100{,}000$, comparing two strategies: \emph{independent grids} ($k = 50$ per quantile, total $7 \times 50 = 350$ thresholds per epoch) and Bonferroni-corrected independent grids ($\alpha' = \alpha/7$ per quantile).

\begin{table}[t]
\caption{Multiple quantile tracking results ($n = 100{,}000$, 5 seeds). Joint coverage requires all 7 quantile CIs to simultaneously contain their true values at all time points. $\uparrow$ means higher is better.}
\label{tab:multiquantile}
\centering
\begin{tabular}{llccc}
\toprule
Distribution & Strategy & Per-quantile Cov.\ $\uparrow$ & Joint Cov.\ $\uparrow$ & Memory \\
\midrule
\multirow{2}{*}{Gaussian} & Independent & 6.9/7 & 3/5 & 17{,}850 \\
                          & Bonferroni  & 7/7   & 5/5 & 17{,}850 \\
\midrule
\multirow{2}{*}{Student-$t$} & Independent & 7/7   & 5/5 & 17{,}850 \\
                             & Bonferroni  & 7/7   & 5/5 & 17{,}850 \\
\bottomrule
\end{tabular}
\end{table}

Table~\ref{tab:multiquantile} shows that independent grids with Bonferroni correction achieve perfect joint coverage (5/5 seeds) on both distributions, at the cost of slightly wider per-quantile CIs. Without Bonferroni correction, joint coverage drops to 3/5 for Gaussian due to failures at $p = 0.9$: the 0.9 quantile falls at $Q_{0.9} = 1.282$, a region where the Gaussian density is relatively flat, making grid placement less precise. The failures occur specifically for seeds 0 and 3, where the adaptive grid in a middle epoch happens to place no threshold sufficiently close to $Q_{0.9}$ (the nearest grid point is more than one grid spacing away). This suggests that for multi-quantile tracking, Bonferroni correction or slightly denser grids at problematic quantile levels are advisable.

On Student-$t$ with independent grids, all 7 quantile CIs maintained simultaneous coverage across all 5 seeds even without Bonferroni correction, because the heavier tails produce wider CIs that provide implicit protection against joint coverage failure.

\subsection{Ablation Studies (Experiment 6)}

We systematically evaluate each design component of \ours{} on Gaussian($0,1$), $p = 0.5$, $k = 50$, $n = 100{,}000$.

\begin{table}[t]
\caption{Ablation studies. Final CI width at $t = 100{,}000$ (mean $\pm$ std, 5 seeds). Each row modifies one component from the default \ours{} configuration (top row). $\downarrow$ means lower is better.}
\label{tab:ablation}
\centering
\begin{tabular}{llccc}
\toprule
Component & Variant & CI Width $\downarrow$ & Coverage & Memory \\
\midrule
\multicolumn{2}{l}{\textbf{\ours{} (default)}} & $\mathbf{0.037 \pm 0.013}$ & 1.00 & 2{,}550 \\
\midrule
\multirow{2}{*}{Grid type} & Fixed grid & $0.165 \pm 0.044$ & 1.00 & 2{,}550 \\
                           & \textit{Improvement} & \multicolumn{3}{c}{\textit{77.3\% narrower with adaptive}} \\
\midrule
\multirow{2}{*}{CS type}   & Hoeffding & $0.037 \pm 0.012$ & 1.00 & 2{,}550 \\
                           & Wald (asymptotic) & $0.000$ & 0.00 & 2{,}550 \\
\midrule
\multirow{2}{*}{Epoch schedule} & Tripling & $0.035 \pm 0.007$ & 1.00 & 1{,}800 \\
                                & Fixed (1000) & $0.095 \pm 0.010$ & 1.00 & 15{,}150 \\
\midrule
\multirow{1}{*}{Intersection} & Without & $0.049 \pm 0.002$ & 1.00 & 2{,}550 \\
\bottomrule
\end{tabular}
\end{table}

\textbf{Adaptive vs.\ fixed grid (Ablation A).} The adaptive grid reduces final CI width by 77.3\% on Gaussian and 90.7\% on Student-$t$ compared to a fixed uniform grid spanning the full data range. The fixed grid's resolution is limited by covering the entire support, while the adaptive grid concentrates thresholds where they matter---around the target quantile. Both maintain valid coverage.

\textbf{Confidence sequence type (Ablation C).} The empirical Bernstein and Hoeffding CSs perform comparably (both achieving coverage $= 1.0$), while the asymptotic Wald interval collapses to zero width with zero coverage---confirming that anytime-valid CSs are essential, not merely preferable. The Wald interval's failure is expected: its $O(1/\sqrt{n})$ width shrinks faster than the LIL rate, violating the conditions for uniform-over-time validity.

\textbf{Epoch schedule (Ablation D).} Doubling ($0.037$) and tripling ($0.035$) epochs perform similarly, while fixed-length epochs ($0.095$) are $2.6\times$ wider and use $6\times$ more memory ($15{,}150$ vs.\ $2{,}550$). The fixed schedule creates too many short epochs, each with too few observations for the Bernoulli CS to be informative, while geometric schedules concentrate data in longer later epochs where CIs are tightest.

\textbf{Multi-epoch intersection (Ablation E).} Intersecting CIs across epochs reduces width by 24.6\% ($0.037$ vs.\ $0.049$), confirming the value of multi-resolution refinement. The improvement comes primarily from early epochs whose wide grids ``anchor'' the CI while later epochs refine it.

\begin{figure}[t]
\centering
\includegraphics[width=0.85\linewidth]{figures/figure5_ablations.pdf}
\caption{Ablation summary showing CI width across design choices. Adaptive grids and geometric epoch schedules are the most impactful components.}
\label{fig:ablation}
\end{figure}

\section{Discussion and Limitations}
\label{sec:discussion}

\textbf{Strengths.} \ours{} is the first method providing anytime-valid quantile CIs with sublinear memory, achieving a practical sweet spot at $k = 50$: CI widths within $1.44\times$ of full-memory baselines with $333\times$ memory savings. The adaptive grid is the most impactful design choice (77--91\% width reduction), and the algorithm is simple to implement ($\sim$150 lines of Python) and deploy.

\textbf{Minimum grid size.} The $k = 5$ failure mode (Section~4.3) reveals a fundamental resolution limit: the grid must be dense enough for at least one threshold to fall near $\Qp$ in early epochs, or the adaptive refinement cannot initialize. The required minimum $k$ depends on the initial range and the distribution's density at $\Qp$: distributions with higher density at the target quantile (e.g., Gaussian near the median) tolerate smaller $k$ than those with lower density (e.g., Pareto at $p = 0.95$). We recommend $k \geq 20$ as a safe minimum and $k = 50$ as a robust default; an adaptive scheme that detects and recovers from grid starvation could further reduce the minimum.

\textbf{Stationarity assumption.} Our method assumes i.i.d.\ observations within each epoch. In nonstationary streams where the distribution shifts, the quantile $\Qp$ itself changes over time, and our CI targets the population quantile of the epoch's distribution rather than tracking drift. Incorporating the importance-weighted CSs of \citet{mineiro2023time} could address this but adds complexity.

\textbf{Baseline comparison.} Our experiments compare against a reservoir-sampling baseline rather than a quantile sketch (e.g., KLL or DDSketch) combined with bootstrap. While a sketch-based quantile estimator would produce tighter point estimates than reservoir sampling, the fundamental limitation we demonstrate---coverage degradation under repeated monitoring---is inherent to all non-sequential CI methods and would affect KLL+bootstrap or DDSketch+bootstrap equally. The reservoir baseline thus provides a favorable comparison for sketch-based approaches, as the full reservoir is available for bootstrap resampling without the information loss from sketch compression.

\textbf{Multi-quantile joint coverage.} When tracking multiple quantile levels simultaneously, joint coverage requires Bonferroni correction ($\alpha' = \alpha/m$ for $m$ quantiles). Without correction, we observed joint coverage failures at the $p = 0.9$ level for Gaussian data (2/5 seeds), where the relatively flat density makes grid placement less reliable. This is a well-understood multiplicity issue, not a flaw in the method, and the Bonferroni correction resolves it completely at the cost of wider per-quantile CIs.

\textbf{Convergence analysis scale.} Our convergence rate analysis uses $n = 10^6$ observations rather than the originally planned $n = 10^7$, covering checkpoints from $t = 100$ to $t = 10^6$. While a longer stream would confirm the convergence exponent more precisely, the 4 orders of magnitude in our data yield a fitted exponent of $-0.516$ (vs.\ theoretical $-0.5$), which is well within the expected range given log-log corrections.

\textbf{Computational cost.} Each observation requires $O(k)$ comparisons against grid thresholds. At $k = 50$, the per-observation cost is negligible---processing 100{,}000 observations takes approximately 0.5 seconds.

\section{Conclusion}
\label{sec:conclusion}

We introduced \ours{}, an algorithm for anytime-valid quantile confidence intervals from streaming data using $O(k \log t)$ memory. By maintaining Bernoulli confidence sequences at an adaptive grid of CDF thresholds organized in doubling epochs, and inverting these CDF bounds, \ours{} bridges the gap between memory-efficient streaming sketches and statistically valid confidence sequences.

Our experiments demonstrate that \ours{} achieves $\geq 0.95$ anytime coverage across four distributions and two quantile levels (100 trials per configuration), with CI widths converging at the optimal $\sqrt{t^{-1}\log\log t}$ rate (exponent $-0.516$). At $k = 50$, it uses $333\times$ less memory than full-memory methods while keeping CIs within $1.44\times$ of their width. On real S\&P~500 data (6{,}288 daily returns), \ours{} produces VaR confidence intervals comparable to full-memory baselines with $3.2\times$ less memory. Unlike bootstrap alternatives, \ours{} maintains valid coverage under continuous monitoring where bootstrap coverage drops to 0.03--0.23.

Future work includes extending \ours{} to nonstationary streams via importance-weighted confidence sequences, developing adaptive grid sizing that detects and recovers from grid starvation (the $k = 5$ failure mode), and optimizing grid placement for extreme quantiles using non-uniform spacing informed by density estimates.

\bibliographystyle{plainnat}
\bibliography{references}

\appendix
\section{Per-Seed Results for Application Scenarios}
\label{app:perseed}

Table~\ref{tab:realworld_seeds} reports per-seed final CI widths for the network latency monitoring experiment (Section~\ref{sec:experiments}, Experiment~4), providing transparency into the variability of results.

\begin{table}[h]
\caption{Per-seed final CI widths for network latency p99 monitoring ($n = 500{,}000$). Seeds are $\{42, 123, 456, 789, 1024\}$. True quantile $= 126.7$.}
\label{tab:realworld_seeds}
\centering
\begin{tabular}{rccc}
\toprule
Seed & \ours{} Width & Full-Memory Width & Bootstrap Width \\
\midrule
42   & 14.52 & 9.07 & 100.39 \\
123  & 12.76 & 9.17 & 69.38 \\
456  & 12.90 & 9.57 & 120.88 \\
789  & 13.43 & 8.51 & 64.75 \\
1024 & 11.74 & 8.30 & 162.14 \\
\midrule
Mean $\pm$ Std & $13.1 \pm 1.0$ & $8.9 \pm 0.5$ & $103.5 \pm 36.4$ \\
\bottomrule
\end{tabular}
\end{table}

Table~\ref{tab:realworld_financial_seeds} reports per-seed results for the financial VaR monitoring experiment using real S\&P~500 daily log returns from Yahoo Finance (January 2000--December 2024, 6{,}288 observations). Seed~0 uses the original chronological ordering; seeds 1--4 use random permutations of the data to assess robustness. Full-sample empirical 5th percentile: $-0.0187$.

\begin{table}[h]
\caption{Per-seed final CI widths for financial VaR monitoring (5th percentile of S\&P~500 daily log returns, $n = 6{,}288$). Real data from Yahoo Finance.}
\label{tab:realworld_financial_seeds}
\centering
\begin{tabular}{rccc}
\toprule
Seed & \ours{} Width & Full-Memory Width & Bootstrap Width \\
\midrule
0 (chrono.) & 0.0000 & 0.0037 & 0.0063 \\
1  & 0.0068 & 0.0039 & 0.0089 \\
2  & 0.0042 & 0.0038 & 0.0090 \\
3  & 0.0054 & 0.0039 & 0.0071 \\
4  & 0.0062 & 0.0038 & 0.0068 \\
\midrule
Mean $\pm$ Std & $0.0045 \pm 0.0024$ & $0.0038 \pm 0.0001$ & $0.0076 \pm 0.0011$ \\
\bottomrule
\end{tabular}
\end{table}

All seeds show \ours{} maintaining valid coverage with CI widths within $1.1$--$1.6\times$ of the full-memory baseline on the latency task, and comparable widths on the financial task. Bootstrap width variability is notably higher (coefficient of variation $35\%$ for latency vs.\ $8\%$ for \ours{}), reflecting the instability of reservoir-based estimation for tail quantiles.

\end{document}
